%% ma: L10-B13-0022
%% lop: 10
%% chuong: 5
%% bai: 13
%% dang: 10.13.4
%% loai: TN4
%% muc: TH
%% dap_an: "C"
%% nguon: "chuyen-de-cac-so-dac-trung-cua-mau-so-lieu-khong-ghep-nhom-toan-10-chuong-trinh-moi.pdf (GV. Phan Nhật Linh) (P02), Bài 02 Dạng 1 Phần I Câu 12"
%% kien_thuc: "Lớp 10 Bài 13 (trung vị)"
%% trang_thai: chinh-thuc
%% ghi_chu: "Tự duyệt theo yêu cầu giáo viên 22/10/2026."
\begin{ex}
Cho bảng tần số về sản lượng cà phê thu được trong $1$ năm (kg/sào) của $20$ hộ gia đình:
\begin{center}
\begin{tabular}{|c|c|c|c|c|c|c|c|}\hline
Sản lượng & $111$ & $112$ & $113$ & $114$ & $115$ & $116$ & $117$\\\hline
Tần số & $1$ & $3$ & $4$ & $5$ & $4$ & $2$ & $1$\\\hline
\end{tabular}
\end{center}
Trung vị của mẫu số liệu trên là
\choice{$117$}{$113{,}5$}{\True $114$}{$111$}
\loigiai{$n=20$ nên $M_e$ là trung bình cộng của giá trị thứ $10$ và thứ $11$. Tần số tích lũy: $1;4;8;13$ nên cả hai giá trị đều là $114$. Vậy $M_e=114$. Chọn C.}
\end{ex}
